\documentclass[11pt]{beamer}
\usetheme{Madrid}
\usepackage{graphicx}
\usepackage{xcolor}
\usepackage{listings}

\lstset{
  basicstyle=\ttfamily\footnotesize,
  breaklines=true,
}

\title[B.Tech Project Evaluation 3]{B.Tech Project Evaluation 3 -- Initial Implementation and Synchronization}
\subtitle{A Rubric \#R3 Evaluation}
\author{Itesh Jung Shahi \and Amit Kumar}
\institute{Associate Professor, Computer Science \& Engineering\\
  Sharda School of Engineering and Technology}
\date{Nov 2025}

\begin{document}

% ------------------------------------------------------------------ title
\begin{frame}
  \titlepage
\end{frame}

% ------------------------------------------------------------------ approval
\begin{frame}{Approval from Guide for the Evaluation}
  \begin{itemize}
    \item APPROVAL FROM GUIDE FOR THE EVALUATION
    \item Paste the screenshot of email
  \end{itemize}
\end{frame}

% ------------------------------------------------------------------ index
\begin{frame}{Index}
  \tableofcontents
\end{frame}

% ------------------------------------------------------------------ introduction
\begin{frame}{Introduction}
  \textbf{INFRA-CVE: A Provenance-Grounded, Confidence-Tiered Automatic CVE Checker for Entire Infrastructure}
  \pause
  Final-year cybersecurity project. Built on the finding that existing infrastructure CVE
  tooling is measured at \textbf{92\% false positives} even on verified SBOMs.
\end{frame}

% ------------------------------------------------------------------ implementation
\begin{frame}{Implementation -- Latest Version}
  \small
  \begin{itemize}
    \item \textbf{Discovery}: banner probing (TCP/TLS/HTTP), host package database (dpkg / apk),
      manifests (package.json / go.mod / requirements.txt)
    \item \textbf{Advisory ingest}: OSV (querybatch -> per-vuln, per-release filtering, ranges,
      severity)
    \item \textbf{Applicability}: distro-first, backport-aware, approximate handling
    \item \textbf{Verification}: corroboration, safe checks, confidence tiers
    \item \textbf{Scoring / reporting}: priority scoring, markdown report, review queue
  \end{itemize}
\end{frame}

% ------------------------------------------------------------------ results 1
\begin{frame}{Results -- Live Target Scan (Banner)}
  \centering
  \includegraphics[width=\linewidth]{/home/kali07/Desktop/Work/College/pbl sem 7/banner_scan.png}
  \pause
  \small
  Live banner scan of \texttt{127.0.0.1:22,80} on Debian 12: two observations, 386 advisory
  records, 386 votes, 58 findings -- all weak (banner-only evidence, no distro revision).
\end{frame}

% ------------------------------------------------------------------ results 2
\begin{frame}{Results -- Host Package Inventory Scan}
  \centering
  \includegraphics[width=\linewidth]{/home/kali07/Desktop/Work/College/pbl sem 7/inventory_scan.png}
  \pause
  \small
  Saved Debian 12 dpkg inventory via \texttt{--inventory} + \texttt{--fixed-only}: 88 packages,
  155 advisory records, 58 findings with a published fix (56 strongly inferred, 2 weak).
\end{frame}

% ------------------------------------------------------------------ results 3
\begin{frame}{Results -- Web Application Scan (Juice Shop)}
  \centering
  \includegraphics[width=\linewidth]{/home/kali07/Desktop/Work/College/pbl sem 7/juice_shop.png}
  \pause
  \small
  Juice Shop source scan: 147 observations, 125 advisories, 71 findings; top critical:
  CVE-2015-9235 (jsonwebtoken 0.4.0 inside \texttt{[0, 4.2.2)}) with fix published.
\end{frame}

% ------------------------------------------------------------------ delivered status
\begin{frame}{Delivered Status}
  \begin{itemize}
    \item Working MVP is ready
    \item Discovery -> CVE check -> verification -> report works end-to-end
    \item Full test suite green (432 tests), type check clean
    \item Honest confidence tiers and provenance kept -- no forced guesses
  \end{itemize}
\end{frame}

% ------------------------------------------------------------------ future implementation
\begin{frame}{Future Implementation}
  \small
  \begin{itemize}
    \item Live authenticated scans of real infrastructure
    \item More ecosystems (Ruby, PHP, Scala, Node runtime-level SBOM correlation)
    \item CI/CD pipeline integration for continuous CVE gating
    \item Container image / SBOM-driven discovery at scale
    \item Issue-tracker / ticketing-system export of findings
  \end{itemize}
\end{frame}

% ------------------------------------------------------------------ conclusion
\begin{frame}{Conclusion}
  \begin{itemize}
    \item MVP delivered and validated
    \item Live demo on Debian 12 and web app (Juice Shop) both produce findings
    \item Honest confidence tiers and provenance -- no forced guesses
  \end{itemize}
\end{frame}

% ------------------------------------------------------------------ outcome
\begin{frame}{Outcome}
  \begin{center}
    INFRA-CVE -- a cybersecurity software that autonomously checks entire
    infrastructure against known CVEs with provenance and honest confidence.
  \end{center}
\end{frame}

\end{document}
